\SourceLeaf{272}$\xi$ ein Object; $\xi'$ dessen Bild; $F$ der erste, $F'$ der zweite Brennpunkt; $E$ der erste, $E'$ der zweite Hauptpunkt; endlich $D$ ein mit der Linse (oder dem Linsensystem) in fester Verbindung stehender Punkt. Mit denselben Buchstaben werden, wie immer, die Coordinaten dieser Punkte in jedem Versuche bezeichnet. Wir setzen ferner die Brennweite $=f$, und die Entfernung des Punktes $D$ von den Brennpunkten, $D-F=p$, $F'-D=q$. Die drei Grössen $f,p,q$ können als die Elemente der Linse betrachtet werden, und zu ihrer Ausmittelung werden also immer drei Versuche erforderlich sein, indem in drei verschiedenen Lagen des Objects und seines Bildes gegen die Linse die Entfernungen derselben von dem Punkte $D$ gemessen werden müssen, welche Aufgabe wir zuvörderst ganz allgemein auflösen wollen.

Die Werthe von $D-\xi$ und $\xi'-D$ seien in einem Versuche $a,b$; in einem zweiten $a',b'$; in einem dritten $a'',b''$. Die allgemeine Gleichung
\[
(F-\xi)(\xi'-F')=ff
\]
gibt uns demnach
\[
\begin{aligned}
(a-p)(b-q)&=ff\\
(a'-p)(b'-q)&=ff\\
(a''-p)(b''-q)&=ff
\end{aligned}
\]
woraus durch Elimination leicht gefunden wird
\[
\begin{aligned}
p&=a-\frac{(a'-a)(a''-a)(b'-b'')}{R}\\
q&=b-\frac{(b-b')(b-b'')(a''-a')}{R}\\
ff&=\frac{(a'-a)(a''-a)(a''-a')(b-b')(b-b'')(b'-b'')}{RR}
\end{aligned}
\]
indem zur Abkürzung
\[
(a''-a)(b-b')-(a'-a)(b-b'')=R
\]
geschrieben wird. Man kann $R$ auch in folgende Form setzen
\[
\begin{aligned}
R&=(a''-a')(b-b')-(a'-a)(b'-b'')\\
&=(a''-a')(b-b'')-(a''-a)(b'-b'')
\end{aligned}
\]
so wie $p$ und $q$ in folgende
\SourceLeaf{273}
\[
\begin{aligned}
p&=a'-\frac{(a'-a)(a''-a')(b-b'')}{R}\\
&=a''-\frac{(a''-a)(a''-a')(b-b')}{R}\\
q&=b'-\frac{(b-b')(b'-b'')(a''-a)}{R}\\
&=b''-\frac{(b-b'')(b'-b'')(a'-a)}{R}
\end{aligned}
\]

\Article{16}
Der allgemeinen im vorhergehenden Artikel gegebenen Auflösung müssen noch einige Bemerkungen beigefügt werden.

I. Es ist vorausgesetzt, dass in den drei Versuchen das Object auf einer und derselben Seite der Linse liegt. Findet man zweckmässig, in einem der Versuche die Linse in verkehrter Lage anzuwenden, so kann man sich denselben so vorstellen, als ob das Bild der Gegenstand und der Gegenstand das Bild wäre, wodurch dieser Fall auf den vorigen zurückgeführt wird.

II. Für sich allein betrachtet, lässt die Formel für $ff$ noch unbestimmt, ob $f$ positiv oder negativ zu nehmen sei: dies entscheidet sich aber schon durch die aufrechte oder verkehrte Stellung des Bildes, indem $\xi'-F'$ und $f$ im ersten Fall entgegengesetzte, im zweiten gleiche Zeichen haben müssen. Auch darf nicht unbemerkt bleiben, dass bei aller Allgemeingültigkeit der analytischen Auflösung doch die praktische Anwendbarkeit auf den Fall wirklicher Bilder (also für einzelne Linsen auf positive Brennweiten) beschränkt bleibt, wenn nicht besondere Hülfsmittel zur Bestimmung des Platzes virtueller Bilder zugezogen werden.

III. Da die Ausführung der Versuche immer nur einen gewissen beschränkten Grad von Schärfe zulässt, so ist es für die Zuverlässigkeit der Resultate keinesweges gleichgültig, was für Combinationen gewählt werden. Im Allgemeinen kann als Regel gelten, dass durch drei Versuche, von denen zwei unter wenig verschiedenen Umständen gemacht sind, jedenfalls nicht alle drei Elemente mit Schärfe bestimmt werden können.

\Article{17}
An einer einfachen Linse sowohl, als an einer solchen, die aus zweien oder mehrern sehr nahe zusammenliegenden zusammengesetzt ist (wie an achromatischen Objectiven von der gewöhnlichen Einrichtung), stehen die beiden Haupt\SourceLeaf{274}punkte in geringer Entfernung von einander. Dürfte man diese Entfernung $E'-E=\lambda$ wie eine bekannte Grösse betrachten, so würden \textit{zwei} Versuche zureichend sein, indem die Gleichung
\[
p+q=2f+\lambda
\]
die Stelle des dritten Versuches vertritt. Verbindet man mit derselben die beiden andern
\[
\begin{aligned}
(a-p)(b-q)&=ff\\
(a'-p)(b'-q)&=ff
\end{aligned}
\]
so erhält man nach der Elimination von $p$ und $q$ zur Bestimmung von $f$ die Gleichung
\[
\begin{aligned}
\frac{(a'+b'-a-b)^2}{(a'-a)(b-b')}.ff+2(a+b+a'+b'-2\lambda)f\\
{}-(a+b'-\lambda)(a'+b-\lambda)=0
\end{aligned}
\]
Diese quadratische Gleichung geht in eine lineare über, wenn $a'+b'-a-b=0$ wird, d. i. wenn die beiden Versuche so angeordnet sind, dass die Entfernung des Bildes vom Objecte in beiden dieselbe bleibt, während die Linse darin zwei verschiedene Stellen einnimmt. Es sei diese Entfernung $=c$, also $a=c-b$, $a'=c-b'$; dadurch wird
\[
4(c-\lambda)f=(c-\lambda+b'-b)(c-\lambda-b'+b)
\]
oder
\[
f=\tfrac14(c-\lambda)-\frac{(b'-b)^2}{4(c-\lambda)}
\]
Für jeden vorgeschriebenen Werth von $c$ muss nemlich $F-\xi$ der Gleichung
\[
F-\xi+\frac{ff}{F-\xi}=F-\xi+\xi'-F'=c-2f-\lambda
\]
Genüge leisten. deren zwei Wurzeln
\[
\begin{aligned}
F-\xi&=\tfrac12(c-2f-\lambda)+\tfrac12\BareRoot{(c-4f-\lambda)(c-\lambda)}\\
F-\xi&=\tfrac12(c-2f-\lambda)-\tfrac12\BareRoot{(c-4f-\lambda)(c-\lambda)}
\end{aligned}
\]
reell und ungleich sind, wenn $c$ grösser ist als $4f+\lambda$, so dass es dann für ein festes Object $\xi$ immer zwei verschiedene Lagen der Linse gibt, bei welchen das Bild mit dem Punkte $\xi+c$ zusammenfällt. Das Product dieser beiden Werthe von $F-\xi$, d. i. $(a-p)(a'-p)$ wird $=ff$, woraus zugleich erhellt,
\SourceLeaf{275}dass $a'-p=b-q$ und $b'-q=a-p$ wird, folglich
\[
\begin{aligned}
p&=\tfrac12(2f+c+\lambda-b-b')\\
q&=\tfrac12(2f-c+\lambda+b+b')\\
E&=D+\tfrac12(b+b'-c)-\tfrac12\lambda\\
E'&=D+\tfrac12(b+b'-c)+\tfrac12\lambda
\end{aligned}
\]

\Article{18}
Bei derjenigen Stellung der Linse, wo $F-\xi=f$ wird, ist $\xi'-\xi=4f+\lambda$, oder das Bild in der kleinsten Entfernung vom Gegenstande; es entfernt sich von demselben, sobald man die Linse aus jener Stellung nach der einen oder nach der andern Seite wegrückt, aber offenbar anfangs sehr langsam. Es folgt daraus, dass wenn für $c$ ein die Grösse $4f+\lambda$ nur wenig überschreitender Werth gewählt ist, die Versuche zur Ausmittelung der beiden erforderlichen Stellungen der Linse oder der Werthe von $b$ und $b'$ nur eine vergleichungsweise geringe Schärfe zulassen. Diese Unsicherheit fällt in ihrer ganzen Stärke auf die Bestimmung von $E$ und $E'$, daher zu \textit{diesem} Zweck die Anwendung des Verfahrens unter solchen Umständen nicht wohl zu gebrauchen ist. Anders aber verhält es sich, wenn es nur darauf ankommt, die Brennweite zu bestimmen, wo die Schärfe durch jenen Umstand Nichts verliert, weil in den Ausdruck für $f$ nur das Quadrat von $b'-b$ eintritt. Die Ausübung des Verfahrens ist überdies in diesem Falle um so bequemer, weil ausser der Distanz $c$ nur die Grösse der Verschiebung der Linse $b'-b$ gemessen zu werden braucht, also die absoluten Werthe von $b$ und $b'$ unnöthig sind.

\Article{19}
Wenn man $\lambda$ ganz vernachlässigt, also
\[
f=\tfrac14c-\frac{(b'-b)^2}{4c}
\]
setzt, so kommt das beschriebene Verfahren mit demjenigen überein, welches \textsc{Bessel} im 17. Bande der Astronomischen Nachrichten vorgeschlagen, und auf die Bestimmung der Brennweite des Objectivs des Königsberger Heliometers angewandt hat. Die strenge Formel zeigt, dass bei der Vernachlässigung von $\lambda$ die Brennweite um
\SourceLeaf{276}
\[
\tfrac14\lambda+\frac{\lambda(b'-b)^2}{4c(c-\lambda)}
\]
zu gross gefunden wird, wo der zweite Theil unter den erwähnten Umständen als unmerklich betrachtet werden kann. Zur Gewinnung eines der Schärfe, welche das Verfahren an sich verstattet, angemessenen Resultats bleibt daher die Berücksichtigung von $\lambda$ wesentlich nothwendig: nur hat es einige Schwierigkeit, sich eine genaue Kenntniss dieser Grösse zu verschaffen. Für eine einfache Linse wird es hinlänglich sein, aus der gemessenen Dicke derselben und dem nothdürftig bekannten Brechungsverhältnisse für $\lambda$ den oben Art. 13 gegebenen Näherungswerth zu berechnen. Auch für eine achromatische Doppellinse mag man allenfalls, in sofern man sich eine genaue Kenntniss von der Dicke jedes einzelnen Bestandtheils verschaffen kann, sich des oben Art. 14 angeführten genäherten Werthes bedienen. Um wenigstens ungefähr eine Vorstellung von dem Einflusse, welchen die Vernachlässigung von $\lambda$ haben kann, zu erhalten, wollen wir, Beispiels halber, ein Objectiv betrachten, wo die Dicke der Kronglaslinse 7 Linien, die Dicke der Flintglaslinse 3 Linien beträgt, und das Brechungsverhältniss für die erstere zu 1,528, für die andere zu 1,618 annehmen. Dadurch wird die Entfernung der beiden Hauptpunkte von einander näherungsweise
\[
\begin{array}{lr}
\text{für die Kronglaslinse}&2{,}42\\
\text{für die Flintglaslinse}&1{,}15\\
\text{und für die zusammengesetzte Linse}&3{,}57\ \text{Linien}
\end{array}
\]
also die Brennweite um 0,89 Linien zu gross. An einem Objective von 8 Fuss Brennweite, dem die vorausgesetzten Dimensionen zukämen, würde also der Fehler etwa $\frac1{1300}$ des Ganzen betragen.

\Article{20}
Wenn man die im vorhergehenden Art. angegebene Bestimmungsart von $\lambda$ nicht anwenden kann, oder sich nicht damit begnügen will, so scheint folgender Weg am zweckmässigsten, um durch unmittelbare Versuche dazu zu gelangen.

Man bestimme den Platz des Bildes eines sehr entfernten Gegenstandes (so gut man kann in der Axe der Linse) relativ gegen den Punkt $D$. In sofern man die Entfernung des Objects als unendlich gross betrachten kann, fällt dieses Bild
\SourceLeaf{277}in $F'$, und der gemessene Abstand $F'-D$ gibt also unmittelbar $q$. Man wiederhole den Versuch, indem man die Linse umkehrt, wo also das Bild in $F$ fallen, und seine Entfernung von $D$ den Werth von $p$ geben wird. Für den dritten Versuch bringe man das Object (auf der Seite von $F$) der Linse möglichst nahe, bestimme die Entfernung des Bildes von diesem Object $=\xi'-\xi$, und zugleich die Entfernung $D-\xi=a''$, und setze $\xi'-D=\xi'-\xi-a''=b''$. Man hat folglich
\[
(p-a'')(q-b'')=ff
\]
oder
\[
\lambda=p+q-2\BareRoot{(p-a'')(q-b'')}
\]
Hat man die Messungen in allen drei Versuchen mit grösster Schärfe ausführen können, so sind dadurch allein schon alle drei Elemente $p,q,f$ hinlänglich genau bestimmt, und es bedarf keiner andern weiter. Wünscht man aber $f$ mit einer noch grössern Schärfe zu erhalten, so hat man jene Versuche nur als eine Vorbereitung zu dem Verfahren des 18. Artikels zu betrachten, die den Werth von $\lambda$ liefert. Um klarer zu übersehen, von welchen Momenten die Schärfe in der so erhaltenen Bestimmung von $\lambda$ hauptsächlich abhängt, setzen wir obige Formel für $\lambda$ in folgende Gestalt
\[
\lambda=a''+b''+\frac{(p-a''-\BareRoot{(p-a'')(q-b'')})^2}{p-a''}
\]
und erwägen, dass $p-a''-\BareRoot{(p-a'')(q-b'')}$ den Abstand des Objects im dritten Versuche vom ersten Hauptpunkte, $p-a''$ hingegen den Abstand jenes Objects vom ersten Brennpunkte vorstellt. Es erhellt daraus, dass unter den Statt habenden Umständen der letzte Theil der Formel für $\lambda$ nur sehr klein wird, und sein berechneter Werth von kleinen Ungenauigkeiten in den Werthen von $p,q,a'',b''$ nur wenig afficirt wird, also die Schärfe der Bestimmung von $\lambda$ hauptsächlich nur von der Schärfe der Messung von $\xi'-\xi=a''+b''$ abhängt.

\Article{21}
In Beziehung auf das im vorhergehenden Artikel angegebene Verfahren verdienen ein Paar Bemerkungen hier noch einen Platz.

I. Zur Ausführung des dritten Versuchs, wo das Bild nur ein virtuelles wird, reichen die sonst anwendbaren Mittel nicht aus: folgende Methode verei\SourceLeaf{278}nigt aber Bequemlichkeit und Schärfe. Auf einer ebenen Fläche wird eine Kreislinie beschrieben, so gross oder wenig grösser als der vorspringende Rand der Fassung des Glases, und der Mittelpunkt dieses Kreises durch zwei zarte Kreuzlinien bezeichnet. Das Glas wird mit der Fassung so auf die Fläche gelegt, dass jene mit der Kreislinie concentrisch ist; dann ein zusammengesetztes an einem festen Stative befindliches und mit einem Fadenkreuze versehenes Mikroskop senkrecht darüber gestellt, und in seiner Hülse so verschoben, bis das Bild der Kreuzlinie genau mit dem Fadenkreuze zusammenfällt; endlich wird das Glas weggenommen, und das Mikroskop durch Verschieben in der Hülse der Ebene genähert, bis das Bild der Kreuzlinie abermals mit dem Fadenkreuze des Mikroskops zusammenfällt. Die leicht auf irgend eine Weise scharf zu messende Grösse der letztern Verschiebung ist die Entfernung des Objects (der Kreuzlinie) von seinem durch die Glaslinie producirten Bilde $=\xi'-\xi$. Den Punkt der Axe der Linse, welcher in der den vorspringenden Rand der Fassung berührenden Ebene liegt, kann man als den festen Punkt $D$ selbst annehmen, in welchem Falle $a''=0$, $b''=\xi'-\xi$ wird, oder, wenn ein anderer Punkt $D$ gewählt war, diesen mit jenem durch leicht sich darbietende Mittel vergleichen, um $a''$ zu finden.

II. Wenn die Entfernung des für den ersten und zweiten Versuch benutzten Gegenstandes, zwar gross, aber doch nicht gross genug ist, um sein Bild als mit dem Brennpunkte ganz zusammenfallend betrachten zu können, so ist eine Reduction nöthig, welche man erhält, indem man das Quadrat der Brennweite mit der Entfernung des Gegenstandes dividirt, und diese Reduction ist von den Abständen des Bildes von dem Punkte $D$ abzuziehen, um die Grössen $q$ und $p$ genau zu erhalten: offenbar ist dazu nur eine grob genäherte Kenntniss der Brennweite und der Entfernung nöthig, insofern letztere sehr gross ist. Indessen kann man diese Reduction eben so leicht durch directe strenge richtige Formeln bestimmen. Es sei für den ersten Versuch $a$ der Werth von $D-\xi$, $b$ der Werth von $\xi'-D$; für den zweiten Versuch hingegen (wo die Linse in verkehrter Stellung angewandt wird) bezeichne man die Entfernung $D-\xi$ mit $b'$, und $\xi'-D$ mit $a'$. Auf diese Weise (die mit der Art. 16, I angegebenen auf Eins hinausläuft) erreichen wir den Vortheil, dass die für die drei Versuche Statt findenden Gleichungen
\SourceLeaf{279}
\[
\begin{aligned}
(a-p)(b-q)&=ff\\
(b'-q)(a'-p)&=ff\\
(a''-p)(b''-q)&=ff
\end{aligned}
\]
gleichlautend mit denen sind, von welchen wir im 15. Art. ausgingen, und also auch die durch Elimination daraus abgeleiteten Formeln ihre Gültigkeit behalten. Ist man bei der Ausführung des zweiten Versuchs so zu Werke gegangen, dass der Ort des Bildes im Raume derselbe ist wie im ersten Versuche (was leicht geschehen kann, obwohl der Erfolg bei der vorausgesetzten grossen Entfernung des Gegenstandes gar nicht merklich abgeändert wird, wenn man es damit nicht ängstlich genommen hat), so wird $a+b=b'+a'$, welche Grösse wir mit $c$ bezeichnen, und die Formeln des 15. Art. erhalten dadurch noch einige Vereinfachung. Es wird nemlich, aus der zweiten Formel für $p$ und der ersten für $q$,
\[
\begin{aligned}
p&=a'-\frac{(a'-a'')(b-b'')}{c-a''-b''}\\
q&=b-\frac{(a'-a'')(b-b'')}{c-a''-b''}
\end{aligned}
\]

III. Wenn man auch das Verfahren des 20. Art. nicht zur vollständigen Bestimmung der Elemente gebrauchen, sondern die schärfste Bestimmung der Brennweite der Methode des 17. Art. vorbehalten will, so bleibt doch jenes zugleich das geeignetste, um die Lage der beiden Hauptpunkte festzusetzen. Es wird nemlich
\[
\begin{aligned}
E&=D+\tfrac12(q-p)-\tfrac12\lambda=D+\tfrac12(b-a')-\tfrac12\lambda\\
E'&=D+\tfrac12(q-p)+\tfrac12\lambda=D+\tfrac12(b-a')+\tfrac12\lambda
\end{aligned}
\]

\Article{22}
Für eine einfache Linse und, allgemein zu reden, auch für ein System von Linsen kann man der Brennweite einen bestimmten Werth und den Haupt- und Brennpunkten bestimmte Plätze nur in sofern beilegen, als von Lichtstrahlen von bestimmter Brechbarkeit die Rede ist; für Strahlen von anderer Brechbarkeit erhalten diese Punkte andere Plätze und die Brennweite einen andern Werth, und das nicht homogene Licht von Gegenständen erleidet daher beim Durchgange durch Gläser eine Farbenzerstreuung. Durch eine Zusammensetzung zweier oder
\SourceLeaf{280}mehrerer Linsen aus verschiedenen Glasarten lässt sich diese Farbenzerstreuung aufheben: zur Vollkommenheit eines achromatischen Objectivs wird aber erforderlich sein, dass Parallelstrahlen sich unabhängig von der Farbe in Einem Punkte vereinigen, und zwar nicht bloss solche, die parallel mit der Axe, sondern auch solche, die geneigt gegen die Axe einfallen, oder mit andern Worten, die verschiedenfarbigen Bilder eines ausgedehnten als unendlich entfernt betrachteten Gegenstandes müssen nicht bloss in Eine Ebene fallen, sondern auch gleiche Grösse haben. Die erste Bedingung beruhet auf der Identität des zweiten Brennpunkts für verschiedenfarbige Strahlen, die zweite auf der Gleichheit der Brennweite, und da diese die Entfernung des zweiten Brennpunkts vom zweiten Hauptpunkte ist, so kann man auch die beiden Bedingungen dadurch ausdrücken, dass beide Punkte zugleich für rothe und violette Strahlen dieselben sein müssen. Ist die erste Bedingung allein erfüllt, so geben die gegen die Axe geneigten Strahlen kein reines Bild; allein eine \textit{sehr} geringe Ungleichheit der Brennweiten für verschiedenfarbige Strahlen wird immer als ganz unschädlich betrachtet werden dürfen.

In der Theorie der achromatischen Objective pflegt man nur die erste Bedingung zu berücksichtigen. Allein bei der gewöhnlichen Construction dieser Objective, wo die beiden Linsen entweder in Berührung oder in einem äusserst geringen Abstande von einander sich befinden, wird die Lage der Hauptpunkte von der ungleichen Brechbarkeit der Lichtstrahlen so wenig afficirt, dass die zweite Bedingung von selbst erfüllt ist, wenn nicht genau, doch so nahe, dass eine merkliche Unvollkommenheit nicht daraus entstehen kann: auch lässt sich, wenn man es der Mühe werth hält, die Dicke der Linsen so berechnen, dass eine genaue Identität des zweiten Hauptpunkts für ungleiche Strahlen Statt findet.

Anders verhält es sich hingegen, wenn die convexe Kronglaslinse von der concaven Flintglaslinse durch einen beträchtlichen Abstand getrennt ist. Es lässt sich leicht zeigen, dass bei solchen Bestimmungen für diesen Abstand und die Brennweiten der einzelnen Linsen, wo der zweite Brennpunkt dieses Linsensystems für verschiedenfarbige Strahlen derselbe bleibt, die Brennweite dieses Systems für die violetten Strahlen nothwendig grösser wird als für die rothen, und dass der Unterschied von derselben Ordnung ist wie derjenige, der (im umgekehrten Sinn) bei einfachen Linsen Statt findet. Dasselbe gilt auch noch, wenn (wie bei den sogenannten dialytischen Fernröhren geschieht) anstatt der zweiten Linse
\SourceLeaf{281}eine Zusammensetzung aus einer Flintglaslinse, und einer Kronglaslinse, in Berührung oder sehr geringem Abstande von einander, angenommen wird. Immer bleibt es unmöglich, auf diese Weise von einem ausgedehnten Objecte ein vollkommen farbenreines Bild hervorzubringen, indem das violette Bild, wenn es in demselben Abstande von dem Linsensysteme liegen soll wie das rothe, nothwendig grösser wird, als das letztere.

Man darf jedoch hieraus keinesweges folgern, dass \textit{Fernröhre} von dieser letztern Einrichtung in Beziehung auf Achromatismus unvollkommener bleiben müssen, als Fernröhre mit achromatischen nach der gewöhnlichen Art construirten und ein völlig farbenreines Bild hervorbringenden Objectiven. Man kann vielmehr gerade umgekehrt behaupten, dass jene bei einer wohlberechneten Anordnung der Oculare dem Auge das farbenreinere Bild zu geben fähig sind.

In der That kann ein vollkommen farbenreines vom Objectiv erzeugtes Bild (möge es ein wirkliches oder virtuelles sein) wegen der Farbenzerstreuung, welche durch die Oculargläser hervorgebracht wird, dem \textit{Auge} nicht vollkommen rein \textit{erscheinen}; man verhütet zwar durch besondere Anordnung der Oculare den sogenannten farbigen Rand, kann aber damit die Längenabweichung nicht aufheben, welche noch durch den Umstand vergrössert wird, dass das menschliche Auge selbst nicht achromatisch ist. Man bewirkt nur, dass die letzten Bilder, rothes und violettes, in einerlei scheinbarer Grösse, nicht aber, dass sie in gleichem Abstande oder \textit{zugleich} deutlich erscheinen.

Die ungleiche Grösse der ersten Bilder, des rothen und violetten, welche bei den dialytischen Objectiven unvermeidlich ist, lässt sich aber durch eine angemessene Einrichtung der Oculare sehr wohl compensiren, so dass der farbige Rand in der Erscheinung eben so gut gehoben wird, wie bei Fernröhren von gewöhnlicher Einrichtung, während die zweite eben berührte Unvollkommenheit auch hier bleibt, so lange das erste rothe und violette Bild in gleicher Entfernung von dem Objective liegen.

Es ist also klar, dass um \textit{im Auge} ein vollkommen farbenreines Bild hervorzubringen, das erste Bild eine gewisse von den Verhältnissen der Oculare und dem Nichtachromatismus des menschlichen Auges abhängende Längenabweichung haben \textit{muss}. Theoretisch betrachtet lässt sich nun allerdings auch ein Objectiv von gewöhnlicher Einrichtung so berechnen, dass eine vorgeschriebene Längenabweichung Statt findet; allein abgesehen von der Schwierigkeit, der ganzen
